\subsection{Ryll-Nardzewski 定理}

在本节中，E 表示域 R 上的一个赋范空间，T 表示 E 上的一个 Hausdorff 局部凸拓扑，对于它 E 的范数是下半连续的。这些假设特别在下列情形下满足：

\begin{enumerate}
\def\labelenumi{\alph{enumi})}
\item
  \(\mathcal{T}\) 是由赋范空间 E 的范数诱导的拓扑。
\item
  \(\mathcal{T}\) 是赋范空间 E 的弱化拓扑 \(\sigma (\mathrm{E},\mathrm{E}^{\prime})\)。
\item
  E 是一个赋范空间 F 的对偶且 \(\mathcal{T} = \sigma (\mathrm{F}',\mathrm{F})\)
\end{enumerate}

\(d\) ) 存在两个赋范空间 \(\mathrm{F}_1\) 与 \(\mathrm{F}_2\)，使得 \(\mathrm{E} = \mathcal{L}(\mathrm{F}_1;\mathrm{F}_2)\) 且 \(\mathcal{T}\) 是简单收敛拓扑。

除非另有明确说明，拓扑概念均指拓扑 \(\mathcal{T}\) 而言。

设 K 是 E 的一个凸子集。假设 K（对于拓扑 T）是紧的，且 K 对于由 E 的范数定义的距离满足第一可数公理。

引理 2. --- 假设 K 至少含两个点。对每个 \(\varepsilon > 0\)，存在 K 的一个划分，分成两个非空子集 \(\mathbf{K}_1\) 与 \(\mathbf{K}_2\)，它们具有下列性质：

\begin{enumerate}
\def\labelenumi{\alph{enumi})}
\item
  \(\mathrm{K}_1\) 凸且紧；
\item
  我们有 \(\| x_1 - x_2\| < \varepsilon\)（对每个 \(x_{1}\) 与每个 \(x_{2}\)，二者均属于 \(\mathbf{K}_2\)）。
\end{enumerate}

设 L 是 K 的所有极点所成之集的闭包。由 Krein-Milman 定理 (II, 第 55 页, 定理 1)，K 是 L 的闭凸包。由于 K 至少含两个点，L 亦然。对每个 \(x \in L\)，设 \(A_{x}\) 是所有 \(y \in L\) 中满足 \(\|x - y\| \leqslant \varepsilon/4\) 者所成的集。根据 K 上的假设，存在 L 的一个可数子集 D 使得 \(L = \bigcup_{x \in D} A_{x}\)。由于范数下半连续，诸集

\(A_{x}\) 中的每一个都是闭的。把 Baire 定理 (GT, IX, § 5, No.~3, 定理 1) 应用于紧空间 L，我们看出存在 D 中的一点 a 与 E 中的一个开子集 U，使得 \(L \cap U\) 非空且包含在 \(A_{a}\) 中。由于 L 至少含两个点，且 E 是 Hausdorff 的，我们可以选取 U 使得 \(L \not\subset U\)。

设 M 是 \(L \cap \complement U\) 的闭凸包。对满足 0 \textless{} t \textless{} 1 的每个实数 t，设 \(M_{t}\) 是所有形如 \(tx_{1} + (1 - t)x_{2}\)（其中 \(x_{1} \in M\)、\(x_{2} \in K\)）的向量所成的集；这是 K 的一个非空紧凸子集。我们将用反证法证明 \(M_{t} \neq K\)。假设 \(M_{t} = K\)；那么 K 的每个极点 x 属于 \(M_{t}\)，从而可以写成形式 \(x = tx_{1} + (1 - t)x_{2}\)（其中 \(x_{1} \in M\)、\(x_{2} \in K\)）。这蕴含 \(x = x_{1} = x_{2}\)，于是 \(x \in M\)。由 Krein-Milman 定理 (II, 第 55 页, 定理 1)，我们有 K = M，即 K 是 \(L \cap \complement U\) 的闭凸包。由 II, 第 56 页 的推论，这蕴含 \(L \subset L \cap \complement U\)，它与关系式 \(L \cap U \neq \varnothing\) 矛盾。

令 \(d = \sup_{x \in \mathbf{K}, y \in \mathbf{K}} \| x - y \|\)，并选取实数 \(t\) 满足 \(0 < t < 1\) 与 \(t < \varepsilon / 4d\)。令 \(\mathbf{K}_1 = \mathbf{M}_t\)，\(\mathbf{K}_2 = \mathbf{K} - \mathbf{M}_t\)。由前面的论证，诸集 \(\mathbf{K}_1\) 与 \(\mathbf{K}_2\) 非空，且 \(\mathbf{K}_1\) 凸而紧。设 \(\mathbf{M}'\) 是 \(\mathbf{L} \cap \mathbf{U}\) 的闭凸包。由于 \(\mathbf{K}\) 是集 \(\mathbf{L} = (\mathbf{L} \cap \complement \mathbf{U}) \cup (\mathbf{L} \cap \mathbf{U})\) 的闭凸包，它也是 \(\mathbf{M} \cup \mathbf{M}'\) 的闭凸包。设 \(x_1\) 与 \(x_2\) 是 \(\mathbf{K}_2\) 中的两点；对 \(i = 1, 2\)，存在 \(y_i \in \mathbf{M}\)、\(z_i \in \mathbf{M}'\) 与实数 \(\alpha_i\)，使得 \(0 \leqslant \alpha_i \leqslant 1\) 且 \(x_i = \alpha_i y_i + (1 - \alpha_i) z_i\)。若 \(\alpha_i \geqslant t\)，则 \(x_i = ty_i + (1 - t) \left\{\frac{\alpha_i - t}{1 - t} y_i + \frac{1 - \alpha_i}{1 - t} z_i\right\}\)；这与 \(x_i \notin \mathbf{M}_t\) 的假设矛盾。于是 \(\alpha_i < t\)（对 \(i = 1, 2\)），从而

\[
\left\| x _ {i} - z _ {i} \right\| = \left\| \alpha_ {i} \left(y _ {i} - z _ {i}\right) \right\| = \alpha_ {i} \left\| y _ {i} - z _ {i} \right\| \leqslant \alpha_ {i} d <   d t <   \varepsilon / 4.
\]

对每个点 \(z\)（属于 \(\mathbf{M}'\)），我们有 \(\| z - a \| \leqslant \varepsilon / 4\)（因为 \(\mathrm{L} \cap \mathrm{U} \subset \mathrm{A}_a\)），于是特别地 \(\| z_i - a \| \leqslant \varepsilon / 4\)。这样

\[
\left\| x _ {1} - x _ {2} \right\| \leqslant \sum_ {i = 1} ^ {2} \left(\left\| x _ {i} - z _ {i} \right\| + \left\| z _ {i} - a \right\|\right) <   \varepsilon .
\]

证明完毕。

引理 3. --- 设 G 是 K 上（对于 T）连续的仿射变换所成的一个群。假设 K 非空，且对所有 x、K 中的 y 与 G 中所有 g 有 \(\|gx - gy\| = \|x - y\|\)。那么存在 K 中的一点，它在 G 下不变。

设 \(\mathfrak{J}\) 是 \(K\) 的非空子集中对于 \(G\) 闭、凸且稳定的那些子集所成的族。若 \((L_{\alpha})_{\alpha \in I}\) 是 \(\mathfrak{J}\) 的元素的一个按包含关系全序的族，则集 \(L = \bigcap_{\alpha \in I} L_{\alpha}\) 属于 \(\mathfrak{J}\)。于是 (S, III, § 3, No.~4, 定理 2)，存在

一个元素 \(\mathbf{L}\)（属于 \(\mathfrak{J}\)），它对于包含关系是极小的。我们将证明 \(\mathbf{L}\) 缩为一点。

我们用反证法论证，假设 L 至少含两个不同的点 \(x_{1}\) 与 \(x_{2}\)，令 \(x = (x_{1} + x_{2})/2\)，\(\varepsilon = \|x_{1} - x_{2}\|/2\)。凸紧集 L 对于由范数定义的距离满足第一可数公理 (GT, IX, § 2, No.~8)。于是我们可以应用引理 2，找到 L 的一个紧凸子集 \(L_{1}\)，它异于 \(\varnothing\) 且异于 L，并具有下列性质：

\begin{enumerate}
\def\labelenumi{(\Alph{enumi})}
\tightlist
\item
  对每两个点 \(y_{1}\) 与 \(y_{2}\)（属于 \(L - L_{1}\)），我们有 \(\|y_{1} - y_{2}\| < \varepsilon\)。
\end{enumerate}

我们将用反证法证明 \(gx \in L_{1}\) 对所有 \(g \in G\) 成立。设 \(g \in G\) 使得 \(gx \in L - L_{1}\)，那么我们有

\[
\left\| g x _ {i} - g x \right\| = \left\| x _ {i} - x \right\| = \left\| x _ {1} - x _ {2} \right\| / 2 = \varepsilon ,
\]

（对 \(i = 1, 2\)）。由性质 (A)，我们有 \(gx_{i} \in \mathbf{L}_{1}\)。由于 \(\mathbf{L}_{1}\) 凸，我们断定 \(gx = (gx_{1} + gx_{2}) / 2\) 属于 \(\mathbf{L}_{1}\)，这与假设矛盾。

设 \(L'\) 是 x 的轨道 Gx 的闭凸包。集 \(L'\) 属于 J。由前面的论证，我们有 \(L' \subset L_1\)，从而 \(L' \subset L\)，\(L' \neq L\)。这与 L 的极小性矛盾，证明完成。

定理 2 (Ryll-Nardzewski). --- 设 E 是一个赋范空间，K 是 E 的一个非空凸子集，它对于弱化拓扑 \(\sigma(E, E')\) 是紧的。设 G 是 K 的等距仿射变换所成的一个群。那么存在 K 中的一点，它在 G 下不变。

对每个 \(g \in \mathbf{G}\)，设 \(\mathbf{K}_g\) 表示所有点 \(x\)（属于 \(\mathbf{K}\)）中满足 \(gx = x\) 者所成的集；设 \(\mathbf{K}\) 赋以弱化拓扑；每个集 \(\mathbf{K}_g\) 在紧空间 \(\mathbf{K}\) 中凸且闭。我们将证明交 \(\bigcap_{g \in \mathbf{G}} \mathbf{K}_g\) 非空；为此，只需

证明集 \(K_{g_{1}} \cap \ldots \cap K_{g_{n}}\) 对 G 中任意 \(g_{1}, \ldots, g_{n}\) 都非空。固定 \(g_{1}, \ldots, g_{n}\)，设 H 是 G 的由 \(\{g_{1}, \ldots, g_{n}\}\) 生成的子群。在 K 中选取一点 a，设 L 表示 a 的轨道 Ha 的闭凸包。设 D 是形如 \(\lambda_{1}h_{1}a + \cdots + \lambda_{m}h_{m}a\) 的元素所成的可数集，其中 \(\lambda_{1}, \ldots, \lambda_{m}\) 是满足 \(\lambda_{1} + \cdots + \lambda_{m} = 1\) 的正有理数，而 \(h_{1}, \ldots, h_{m}\) 是 H 中的元素。D 对于强拓扑的闭包 \(\overline{D}\) 是凸的，从而它对于 \(\sigma(\mathrm{E}, \mathrm{E}')\) 是闭的 (IV, 第 4 页, 命题 2)；因此 \(\overline{D} = L\)，这就证明了 L 是一个对于距离 \((x, y) \mapsto \|x - y\|\) 满足第一可数公理的度量空间。现在我们可以应用引理 2。存在 L 中的一点 b，它在 H 下不变，从而 \(b \in K_{g_{1}} \cap \ldots \cap K_{g_{n}}\)。

推论. --- 设 E 是一个自反 Banach 空间，G 是赋范空间 E 的自同构所成的一个群，K 是 E 的一个子集。假设 K 非空、凸、闭、有界且在 G 下稳定。那么存在 K 中的一点，它在 G 下不变。

由于 E 自反，K 对于 \(\sigma(\mathrm{E}, \mathrm{E}')\) 是紧的 (IV, 第 15 页, 定理 1)。此外，G 的每个元素属于 \(\mathcal{L}(\mathrm{E})\)。

